% !TEX program = xelatex
% =====================================================================
% LLM Theory Seminar -- English Study Notes
% Compile with:
%   xelatex FILENAME.tex
%   xelatex FILENAME.tex
% or: latexmk -xelatex FILENAME.tex
% =====================================================================
\documentclass[11pt,a4paper,oneside,openany]{memoir}
\usepackage{amsmath,amssymb,amsthm,mathtools,bm}
\usepackage{booktabs,longtable,array,tabularx,makecell}
\usepackage{enumitem}
\usepackage[most]{tcolorbox}
\usepackage[dvipsnames]{xcolor}
\usepackage[unicode,bookmarks=false]{hyperref}
\usepackage{cleveref}
\usepackage{needspace}
\setlrmarginsandblock{27mm}{27mm}{*}
\setulmarginsandblock{24mm}{28mm}{*}
\checkandfixthelayout
\setlength{\parindent}{0pt}
\setlength{\parskip}{0.48em}
\setlength{\emergencystretch}{3em}
\hbadness=10000
\setlength{\headheight}{15pt}
\linespread{1.055}\selectfont
\setlist[itemize]{topsep=0.28em,itemsep=0.12em,parsep=0pt,leftmargin=1.55em}
\setlist[enumerate]{topsep=0.28em,itemsep=0.16em,parsep=0pt,leftmargin=1.75em}
\definecolor{NoteBlue}{HTML}{294E6B}
\definecolor{NoteBlueLight}{HTML}{F2F6F9}
\definecolor{NoteGray}{HTML}{5A6268}
\definecolor{NoteGrayLight}{HTML}{F6F6F4}
\definecolor{NoteWarm}{HTML}{7A6545}
\definecolor{NoteWarmLight}{HTML}{FAF7F0}
\hypersetup{colorlinks=true,linkcolor=NoteBlue,citecolor=NoteBlue,urlcolor=NoteBlue}
\makepagestyle{seminar}
\makeoddhead{seminar}{\footnotesize LLM Theory Seminar}{}{\footnotesize Transformer Expressivity}
\makeoddfoot{seminar}{}{\footnotesize\thepage}{}
\makeheadrule{seminar}{\textwidth}{0.35pt}
\pagestyle{seminar}
\aliaspagestyle{chapter}{seminar}
\makechapterstyle{seminarnotes}{
  \setlength{\beforechapskip}{8pt}
  \setlength{\midchapskip}{8pt}
  \setlength{\afterchapskip}{22pt}
  \renewcommand{\chapterheadstart}{\vspace*{\beforechapskip}}
  \renewcommand{\chapnamefont}{\normalfont\small\bfseries}
  \renewcommand{\chapnumfont}{\normalfont\bfseries\fontsize{34}{38}\selectfont\color{NoteBlue}}
  \renewcommand{\chaptitlefont}{\normalfont\huge\bfseries\color{black}}
  \renewcommand{\printchaptername}{}
  \renewcommand{\chapternamenum}{}
  \renewcommand{\printchapternum}{\chapnumfont\thechapter}
  \renewcommand{\afterchapternum}{\par\nobreak\color{black}\vskip\midchapskip}
  \renewcommand{\printchaptertitle}[1]{\chaptitlefont ##1}
  \renewcommand{\afterchaptertitle}{\par\nobreak\vskip\afterchapskip}
}
\chapterstyle{seminarnotes}
\setsecheadstyle{\Large\bfseries}
\setsubsecheadstyle{\large\bfseries}
\setsubsubsecheadstyle{\normalsize\bfseries}
\setbeforesecskip{-1.3\baselineskip plus -0.2\baselineskip minus -0.1\baselineskip}
\setaftersecskip{0.55\baselineskip}
\setbeforesubsecskip{-0.9\baselineskip plus -0.15\baselineskip minus -0.1\baselineskip}
\setaftersubsecskip{0.35\baselineskip}
\newcommand{\R}{\mathbb{R}}
\newcommand{\E}{\mathbb{E}}
\newcommand{\norm}[1]{\left\lVert #1\right\rVert}
\newcommand{\inner}[2]{\left\langle #1,#2\right\rangle}
\newcommand{\grad}{\nabla}
\DeclareMathOperator*{\argmin}{arg\,min}
\DeclareMathOperator*{\argmax}{arg\,max}
\DeclareMathOperator{\rank}{rank}
\DeclareMathOperator{\Tr}{Tr}
\tcbset{seminarbox/.style={enhanced,breakable,sharp corners,boxrule=0pt,frame hidden,left=3.2mm,right=3mm,top=1.8mm,bottom=1.8mm,before={\par\Needspace{5\baselineskip}},before skip=0.8em,after skip=0.8em,fonttitle=\bfseries\small,coltitle=black,attach title to upper={\quad}}}
\newtcolorbox{keybox}[1][]{seminarbox,colback=NoteBlueLight,borderline west={1.8pt}{0pt}{NoteBlue},title={Key point#1}}
\newtcolorbox{paperbox}[1][]{seminarbox,colback=NoteGrayLight,borderline west={1.8pt}{0pt}{NoteGray},title={From the original paper#1}}
\newtcolorbox{suppbox}[1][]{seminarbox,colback=NoteWarmLight,borderline west={1.8pt}{0pt}{NoteWarm},title={Supplementary derivation#1}}
\newtcolorbox{warningbox}[1][]{seminarbox,colback=NoteGrayLight,borderline west={1.8pt}{0pt}{NoteGray},title={Caution#1}}
\newtcolorbox{presentbox}{seminarbox,colback=white,borderline west={2.2pt}{0pt}{black!70},fontupper=\footnotesize,top=1.2mm,bottom=1.2mm,before={\par},before skip=0.5em,after skip=0.5em,title={How to say it in the presentation}}
\theoremstyle{definition}
\newtheorem{definition}{Definition}[chapter]
\newtheorem{proposition}{Proposition}[chapter]
\newtheorem{theorem}{Theorem}[chapter]
\begin{document}
\begin{titlingpage}
\thispagestyle{empty}
\vspace*{1.2cm}
{\small\bfseries LLM Theory Seminar \hfill Week 4\par}
\vspace{1.4cm}
{\HUGE\bfseries Transformer Expressivity\par}
\vspace{0.55cm}
{\large English seminar study notes\par}
\vspace{0.9cm}
\rule{\textwidth}{0.7pt}
\vspace{1.1cm}
\begin{minipage}{0.86\textwidth}
\small
These notes are an English companion to the seminar handout.  The emphasis is on
mathematical derivations that can be followed line by line, on separating theorem
statements from interpretation, and on identifying the assumptions under which each
claim is valid.
\end{minipage}
\vfill
{\small\textbf{Primary readings}\\[0.25em]
Yun et al., \textbf{Are Transformers Universal Approximators of Sequence-to-Sequence Functions?}, ICLR 2020.\\Bhojanapalli et al., \textbf{Low-Rank Bottleneck in Multi-head Attention Models}, ICML 2020.
\par}
\vspace{0.8cm}
{\small September 2026\par}
\end{titlingpage}
\tableofcontents*
\clearpage

\chapter{Notation and the three levels of expressivity}
\section{Matrix conventions}
Let a length-$n$ sequence with token dimension $d$ be represented as $X=[x_1,\ldots,x_n]\in\R^{d\times n}$, so columns are tokens.  A single attention head forms
\[
Q=W_QX,\qquad K=W_KX,\qquad V=W_VX,
\]
and outputs, schematically,
\[
\operatorname{Attn}(X)=V\,\operatorname{softmax}(K^\top Q),
\]
where the softmax is column-wise under this convention.  Other papers transpose every matrix; the mathematical content is unchanged if one is consistent.

\section{Three questions that should not be mixed}
\begin{enumerate}
\item \textbf{Block-level representation:} which attention matrices can one head produce?
\item \textbf{Network-level approximation:} which continuous sequence-to-sequence functions can a deep Transformer approximate?
\item \textbf{Efficient representation:} how much depth, width, and precision are required?
\end{enumerate}
Universal approximation answers the second question, not automatically the first or third.

\begin{presentbox}
A Transformer can be universal even if an individual attention head has a rank bottleneck.  Universality is a statement about the composition of many blocks, not about unrestricted one-layer attention matrices.
\end{presentbox}

\chapter{Permutation symmetry of self-attention}
\section{Self-attention without positional information}
Let $P\in\R^{n\times n}$ be a permutation matrix.  Reordering the tokens gives $XP$.  Then
\[
Q(XP)=W_QXP=Q(X)P,
\]
and similarly $K(XP)=K(X)P$, $V(XP)=V(X)P$.  Hence
\[
K(XP)^\top Q(XP)=P^\top K(X)^\top Q(X)P.
\]
Because column-wise softmax commutes with simultaneous row/column permutation,
\[
\operatorname{softmax}(P^\top AP)=P^\top\operatorname{softmax}(A)P.
\]
Therefore
\begin{align*}
\operatorname{Attn}(XP)
&=V(X)P\,P^\top\operatorname{softmax}(K^\top Q)P\\
&=\operatorname{Attn}(X)P.
\end{align*}
A token-wise feed-forward network also commutes with $P$, so the whole Transformer without positional information is permutation equivariant:
\[
\boxed{T(XP)=T(X)P.}
\]

\section{What the symmetry excludes}
If a target function $F$ distinguishes absolute token positions, then $F(XP)\neq F(X)P$ for some $P$.  Such a function cannot be uniformly represented by a position-free Transformer, no matter how it is trained.

\begin{presentbox}
Before discussing universal approximation, identify the symmetry class.  Without positional information, the correct target class is permutation-equivariant sequence-to-sequence functions.
\end{presentbox}

\chapter{Universal approximation: what the theorem actually says}
\section{$L^p$ approximation on a compact domain}
Let $K\subset\R^{d\times n}$ be compact and let $F:K\to\R^{d\times n}$ be continuous.  A typical approximation statement uses
\[
d_p(F,G)=\left(\int_K\|F(X)-G(X)\|_p^p\,dX\right)^{1/p}.
\]
Universal approximation means: for every $\varepsilon>0$, there exists a Transformer $G$ in the specified architectural class such that $d_p(F,G)<\varepsilon$.

\section{Yun et al.: fixed-width universality}
A striking result of Yun et al. is that a Transformer architecture with a constant number of heads and constant head/feed-forward width can approximate every continuous permutation-equivariant sequence-to-sequence function on a compact domain.  In their notation, the construction can be placed in a class often summarized as
\[
\mathcal T_{2,1,4},
\]
corresponding to two heads, scalar head size, and small feed-forward width, while depth is allowed to grow.

\begin{warningbox}
``Fixed width'' does not mean ``constant total complexity.''  The proof can require very large depth as the desired accuracy and input dimension change.
\end{warningbox}

\section{Why the result is plausible}
The architecture need not store an arbitrary high-dimensional function in a single layer.  Depth allows it to build a discrete contextual code, route information conditionally, and then decode the desired output.

\begin{presentbox}
The surprising part of the theorem is not that a large Transformer is expressive.  It is that width can remain bounded while depth carries the burden of universality.
\end{presentbox}

\chapter{The proof idea: contextual mappings}
\section{Why token-wise MLPs are insufficient}
A token-wise map $g$ has the form
\[
G(X)=[g(x_1),\ldots,g(x_n)].
\]
It cannot distinguish the same token value appearing in different sequence contexts.  Attention is needed to make a token representation depend on the rest of the sequence.

\section{Contextual IDs}
The proof constructs an intermediate map in which each token obtains a representation that uniquely identifies both its own quantized value and the relevant global context.  Once these contextual IDs are distinct, a token-wise network can map them to the desired output values.

\section{Three proof stages}
The high-level proof is:
\begin{enumerate}
\item \textbf{Quantize.} Approximate the continuous target by a piecewise-constant function on a fine grid.
\item \textbf{Contextualize and route.} Use an idealized Transformer with hard selections to assign distinguishable contextual codes and implement the grid function.
\item \textbf{Approximate the idealized operations.} Replace hardmax/piecewise gadgets by ordinary softmax and ReLU blocks with sufficiently sharp parameters.
\end{enumerate}
Continuity and compactness control the approximation error introduced by the grid and by replacing idealized operations.

\begin{suppbox}
The logical structure is the same as many classical UAT proofs: first reduce an arbitrary continuous function to a finite combinatorial object, then prove that the architecture can realize that object, then approximate the idealized components by the actual activation family.
\end{suppbox}

\begin{presentbox}
The core role of attention in the proof is to create context-dependent identifiers.  Once a token knows ``who it is in this sequence,'' a small feed-forward network can decode the desired output.
\end{presentbox}

\chapter{Positional encoding as symmetry breaking}
\section{Adding absolute position}
Let $E=[e_1,\ldots,e_n]$ be a positional code and feed $X+E$ into the Transformer.  A token permutation changes $X$ but not the association between $e_j$ and position $j$, so the permutation-equivariance constraint is broken.

\section{Universality for arbitrary fixed-length sequence maps}
Yun et al. show that with trainable positional encodings, the same basic architecture can approximate arbitrary continuous fixed-length sequence-to-sequence functions, not only permutation-equivariant ones.

\section{The exact scope}
The theorem is for a fixed sequence length $n$ and compact input domain.  It does not imply length generalization to $n'>n$, nor does it state that a learned positional encoding extrapolates algorithmically.

\begin{presentbox}
Positional encoding is mathematically a symmetry-breaking device.  It enlarges the target function class from permutation-equivariant maps to arbitrary continuous maps on a fixed-length sequence space.
\end{presentbox}

\chapter{The expressivity of a single attention matrix}
\section{The softmax inverse problem}
Suppose we want a single head to produce a positive stochastic matrix $A$.  Before softmax, the score matrix has the form
\[
S=K^\top Q=X^\top W_K^\top W_QX.
\]
If the projected key/query dimension is $d_k$, then
\[
\rank(S)\le d_k.
\]
Softmax is nonlinear, but this low-dimensional score factorization still constrains which attention patterns can be produced.

\section{Bhojanapalli et al.: sufficient dimension}
Under suitable full-rank assumptions on the input representation and sufficiently large key/query dimension, one can solve for projection matrices whose scores realize a desired positive stochastic attention pattern after softmax.  The construction effectively works by taking a logarithm of the desired probabilities up to column-wise additive constants and factorizing the resulting score matrix.

\section{Low-rank bottleneck}
When $d_k$ is too small relative to sequence length, there exist attention matrices that cannot be represented by a single head on a fixed input.  The obstruction comes from the low-rank score factorization.

\section{Why this does not contradict UAT}
Deep Transformers can use multiple heads, intermediate nonlinear transformations, and many layers.  A restriction on one attention matrix therefore does not imply a corresponding restriction on the function class of the whole deep network.

\begin{presentbox}
There are two different notions of expressivity: the set of attention matrices one layer can instantiate, and the set of functions a deep Transformer can approximate.  The first can be low-rank constrained while the second is universal.
\end{presentbox}

\chapter{Depth, width, and efficiency}
\section{Small width can hide enormous depth}
A universal approximation theorem is existential.  If the construction needs a grid with resolution $\delta$ in a $dn$-dimensional domain, the number of grid cells can scale like
\[
O(\delta^{-dn}),
\]
which exhibits the usual curse of dimensionality.  A fixed-width network may therefore require impractically large depth.

\section{Depth as sequential computation}
With width fixed, depth provides repeated opportunities to compare, route, and rewrite token representations.  In the UAT proof, this sequential resource is what compensates for small head and MLP widths.

\section{What the theorem does not imply}
Universal approximation does not imply:
\begin{itemize}
\item that SGD finds the constructive solution;
\item polynomial sample complexity;
\item efficient depth or parameter count;
\item length extrapolation;
\item computational tractability of inference;
\item interpretability of learned attention weights.
\end{itemize}

\begin{presentbox}
Expressivity is a feasibility statement.  Learnability, efficiency, and generalization are separate questions, and later weeks study exactly those gaps.
\end{presentbox}

\chapter{Presentation checklist}
\section{Eight formulas to keep}
\begin{align*}
&Q=W_QX,\quad K=W_KX,\quad V=W_VX,\\
&T(XP)=T(X)P\quad\text{(no positional information)},\\
&d_p(F,G)=\left(\int_K\|F(X)-G(X)\|_p^p dX\right)^{1/p},\\
&\mathcal T_{2,1,4}\quad\text{(fixed-width UAT construction)},\\
&X\mapsto X+E\quad\text{breaks permutation symmetry},\\
&S=K^\top Q,\\
&\rank(S)\le d_k,\\
&\text{single-head matrix expressivity}\neq\text{deep-network universality}.
\end{align*}

\section{Common mistakes}
\begin{itemize}
\item Omitting the fixed-length and compact-domain assumptions.
\item Saying that positional encoding is required for permutation-equivariant universality.
\item Treating fixed width as an efficiency theorem.
\item Using the rank of $K^\top Q$ as if it directly bounded the rank of the post-softmax matrix.
\item Inferring SGD learnability from an existential construction.
\end{itemize}

\begin{presentbox}
Thirty-second conclusion: self-attention without positional information is permutation equivariant.  Despite strong one-layer structural constraints, sufficiently deep fixed-width Transformers are universal for continuous equivariant maps; positional encoding breaks the symmetry and extends the result to arbitrary continuous fixed-length sequence maps.
\end{presentbox}

\chapter*{Bibliography}
\addcontentsline{toc}{chapter}{Bibliography}
\begin{itemize}
\item Yun, C. et al. \emph{Are Transformers Universal Approximators of Sequence-to-Sequence Functions?} ICLR, 2020.
\item Bhojanapalli, S. et al. \emph{Low-Rank Bottleneck in Multi-head Attention Models}. ICML, 2020.
\end{itemize}
\end{document}
